\errorcontextlines=100
%\documentclass[final]{ltugboat}
\documentclass[twoside,twocolumn,american]{scrartcl}
%% $Id: doublepage2s2c.tex 92 2021-05-18 09:46:55Z herbert $

\input{preamble.ltx}
\usepackage{xcolor,marginnote}

\usepackage{marginnote,showframe}
\setcounter{tocdepth}{2}

\begin{document}
\title{Examples for doublepage floats \newline with bind correction}
\author{Herbert Voß}
\maketitle

\tableofcontents


\onecolumn
\listoffigures

\newpage\null
\twocolumn
\part{Argument \texttt{doubleFULLPAGE}}
\section{Definition on an odd page}

\Blindtext\hvblindtext

\hvblindtext

\subsection{The default -- in the first column}

\begin{lstlisting}
\hvFloat[doubleFULLPAGE]%
  {figure}%
  {\includegraphics[doubleFULLPAGE]{images/rheinsberg}}%
  [A doublepage image with a caption on the right side of the right part.]%
  {A caption for a double-sided image that will be placed on the right side of the
   right-hand part of the illustration. The illustration begins on the left edge of 
   the paper. A short form is used for the LOF. 
   The parameter is \texttt{doubleFULLPAGE}}%
  {fig:doubleFULLPAGE0}
\end{lstlisting}


\marginnote{Fig. \ref{fig:doubleFULLPAGE0}}
\hvFloat[doubleFULLPAGE]%
  {figure}%
  {\includegraphics[doubleFULLPAGE]{images/rheinsberg}}%
  [A doublepage image with a caption on the right side of the right part.]%
  {A caption for a double-sided image that will be placed on the right side of the
   right-hand part of the illustration. The illustration begins on the left edge of 
   the paper. A short form is used for the LOF. 
   The parameter is \texttt{doubleFULLPAGE}}%
  {fig:doubleFULLPAGE0}

\Blindtext

\Blindtext

%\Blindtext

\hvblindtext

\hvblindtext
\hvblindtext

\subsection{The default -- in the second column}


\begin{lstlisting}
\hvFloat[doubleFULLPAGE]%
  {figure}%
  {\includegraphics[doubleFULLPAGE]{images/rheinsberg}}%
  [A doublepage image with a caption on the right side of the right part.]%
  {A caption for a double-sided image that will be placed on the right side of the
   right-hand part of the illustration. The illustration begins on the left edge of 
   the paper. A short form is used for the LOF. 
   The parameter is \texttt{doubleFULLPAGE}}%
  {fig:doubleFULLPAGE02nd}
\end{lstlisting}


\marginnote{Fig. \ref{fig:doubleFULLPAGE02nd}}
\hvFloat[doubleFULLPAGE]%
  {figure}%
  {\includegraphics[doubleFULLPAGE]{images/rheinsberg}}%
  [A doublepage image with a caption on the right side of the right part.]%
  {A caption for a double-sided image that will be placed on the right side of the
   right-hand part of the illustration. The illustration begins on the left edge of 
   the paper. A short form is used for the LOF. 
   The parameter is \texttt{doubleFULLPAGE}}%
  {fig:doubleFULLPAGE02nd}

\Blindtext

%\Blindtext

%\Blindtext

\hvblindtext

\hvblindtext
\hvblindtext

\subsection{\texttt{bindCorr=1cm} -- first column}

\begin{lstlisting}
\hvFloat[doubleFULLPAGE,capWidth=n,capPos=right,bindCorr=1cm]%
  {figure}%
  {\includegraphics[doubleFULLPAGEbindCorr]{images/rheinsberg}}%
  [A doublepage image with a caption on the right side of the right part.]%
  {A caption for a double-sided image that will be placed on the right side of the
   right-hand part of the illustration. The illustration begins on the left edge of 
   the paper. A short form is used for the LOF. 
   The parameter is \texttt{doubleFULLPAGE}}%
  {fig:doubleFULLPAGE1}
\end{lstlisting}

\marginnote{Fig. \ref{fig:doubleFULLPAGE1}}
\hvFloat[doubleFULLPAGE,capWidth=n,capPos=right,bindCorr=1cm]%
  {figure}%
  {\includegraphics[doubleFULLPAGEbindCorr]{images/rheinsberg}}%
  [A doublepage image with a caption on the right side of the right part.]%
  {A caption for a double-sided image that will be placed on the right side of the
   right-hand part of the illustration. The illustration begins on the left edge of 
   the paper. A short form is used for the LOF. 
   The parameter is \texttt{doubleFULLPAGE}}%
  {fig:doubleFULLPAGE1}

\hvblindtext

%\Blindtext

\Blindtext

\Blindtext

%\Blindtext

\hvblindtext

\hvblindtext

\hvblindtext

\subsection{\texttt{bindCorr=1cm} -- second column}

\begin{lstlisting}
\hvFloat[doubleFULLPAGE,capWidth=n,capPos=right,bindCorr=1cm]%
  {figure}%
  {\includegraphics[doubleFULLPAGEbindCorr]{images/rheinsberg}}%
  [A doublepage image with a caption on the right side of the right part.]%
  {A caption for a double-sided image that will be placed on the right side of the
   right-hand part of the illustration. The illustration begins on the left edge of 
   the paper. A short form is used for the LOF. 
   The parameter is \texttt{doubleFULLPAGE}}%
  {fig:doubleFULLPAGE12nd}
\end{lstlisting}

\marginnote{Fig. \ref{fig:doubleFULLPAGE12nd}}
\hvFloat[doubleFULLPAGE,capWidth=n,capPos=right,bindCorr=1cm]%
  {figure}%
  {\includegraphics[doubleFULLPAGEbindCorr]{images/rheinsberg}}%
  [A doublepage image with a caption on the right side of the right part.]%
  {A caption for a double-sided image that will be placed on the right side of the
   right-hand part of the illustration. The illustration begins on the left edge of 
   the paper. A short form is used for the LOF. 
   The parameter is \texttt{doubleFULLPAGE}}%
  {fig:doubleFULLPAGE12nd}

\hvblindtext

%\Blindtext

\Blindtext

%\Blindtext

\hvblindtext
\hvblindtext

\hvblindtext



\subsection{\texttt{bindCorr=3mm} -- first column}
\begin{lstlisting}
\hvFloat[doubleFULLPAGE,capWidth=n,capPos=right,bindCorr=3mm]%
  {figure}%
  {\includegraphics[doubleFULLPAGEbindCorr]{images/rheinsberg}}%
  [A doublepage image with a caption on the right side of the right part.]%
  {A caption for a double-sided image that will be placed on the right side of the
   right-hand part of the illustration. The illustration begins on the left edge of 
   the paper. A short form is used for the LOF. 
   The parameter is \texttt{doubleFULLPAGE}}%
  {fig:doubleFULLPAGE2}
\end{lstlisting}

\marginnote{Fig. \ref{fig:doubleFULLPAGE2}}
\hvFloat[doubleFULLPAGE,capWidth=n,capPos=right,bindCorr=3mm]%
  {figure}%
  {\includegraphics[doubleFULLPAGEbindCorr]{images/rheinsberg}}%
  [A doublepage image with a caption on the right side of the right part.]%
  {A caption for a double-sided image that will be placed on the right side of the
   right-hand part of the illustration. The illustration begins on the left edge of 
   the paper. A short form is used for the LOF. 
   The parameter is \texttt{doubleFULLPAGE}}%
  {fig:doubleFULLPAGE2}


\Blindtext

\Blindtext

%\Blindtext
%\Blindtext

\hvblindtext

\hvblindtext
\hvblindtext




\subsection{\texttt{bindCorr=3mm} -- second column}
\begin{lstlisting}
\hvFloat[doubleFULLPAGE,capWidth=n,capPos=right,bindCorr=3mm]%
  {figure}%
  {\includegraphics[doubleFULLPAGEbindCorr]{images/rheinsberg}}%
  [A doublepage image with a caption on the right side of the right part.]%
  {A caption for a double-sided image that will be placed on the right side of the
   right-hand part of the illustration. The illustration begins on the left edge of 
   the paper. A short form is used for the LOF. 
   The parameter is \texttt{doubleFULLPAGE}}%
  {fig:doubleFULLPAGE22nd}
\end{lstlisting}

\marginnote{Fig. \ref{fig:doubleFULLPAGE22nd}}
\hvFloat[doubleFULLPAGE,capWidth=n,capPos=right,bindCorr=3mm]%
  {figure}%
  {\includegraphics[doubleFULLPAGEbindCorr]{images/rheinsberg}}%
  [A doublepage image with a caption on the right side of the right part.]%
  {A caption for a double-sided image that will be placed on the right side of the
   right-hand part of the illustration. The illustration begins on the left edge of 
   the paper. A short form is used for the LOF. 
   The parameter is \texttt{doubleFULLPAGE}}%
  {fig:doubleFULLPAGE22nd}


\Blindtext

%\Blindtext

%\Blindtext

\hvblindtext
\hvblindtext

\hvblindtext


\subsection{\texttt{bindCorr=<inside textwidth>} -- first column}

\begin{lstlisting}
\hvFloat[doubleFULLPAGE,capWidth=n,bindCorr=\the\dimexpr1in+\oddsidemargin]%
  {figure}%
  {\includegraphics[doubleFULLPAGEbindCorr]{images/rheinsberg}}%
  [A doublepage image with a caption on the right side of the right part.]%
  {A caption for a double-sided image that will be placed on the right side of the
   right-hand part of the illustration. The illustration begins on the left edge of 
   the paper. A short form is used for the LOF. 
   The parameter is \texttt{doubleFULLPAGE}}%
  {fig:doubleFULLPAGE3}
\end{lstlisting}

\marginnote{Fig. \ref{fig:doubleFULLPAGE3}}
\hvFloat[doubleFULLPAGE,capWidth=n,bindCorr=\the\dimexpr1in+\oddsidemargin]%
  {figure}%
  {\includegraphics[doubleFULLPAGEbindCorr]{images/rheinsberg}}%
  [A doublepage image with a caption on the right side of the right part.]%
  {A caption for a double-sided image that will be placed on the right side of the
   right-hand part of the illustration. The illustration begins on the left edge of 
   the paper. A short form is used for the LOF. 
   The parameter is \texttt{doubleFULLPAGE}}%
  {fig:doubleFULLPAGE3}


\Blindtext

\Blindtext

%\Blindtext

\hvblindtext

\hvblindtext
\hvblindtext

\hvblindtext

\subsection{\texttt{bindCorr=<inside textwidth>} -- second column}

\begin{lstlisting}
\hvFloat[doubleFULLPAGE,capWidth=n,bindCorr=\the\dimexpr1in+\oddsidemargin]%
  {figure}%
  {\includegraphics[doubleFULLPAGEbindCorr]{images/rheinsberg}}%
  [A doublepage image with a caption on the right side of the right part.]%
  {A caption for a double-sided image that will be placed on the right side of the
   right-hand part of the illustration. The illustration begins on the left edge of 
   the paper. A short form is used for the LOF. 
   The parameter is \texttt{doubleFULLPAGE}}%
  {fig:doubleFULLPAGE32nd}
\end{lstlisting}

\marginnote{Fig. \ref{fig:doubleFULLPAGE32nd}}
\hvFloat[doubleFULLPAGE,capWidth=n,bindCorr=\the\dimexpr1in+\oddsidemargin,capPos=after]%
  {figure}%
  {\includegraphics[doubleFULLPAGEbindCorr]{images/rheinsberg}}%
  [A doublepage image with a caption on the right side of the right part.]%
  {A caption for a double-sided image that will be placed on the right side of the
   right-hand part of the illustration. The illustration begins on the left edge of 
   the paper. A short form is used for the LOF. 
   The parameter is \texttt{doubleFULLPAGE}}%
  {fig:doubleFULLPAGE32nd}


%\Blindtext

\hvblindtext

\hvblindtext

%\hvblindtext




\section{Definition on an even page}


\subsection{The default -- first column}
\begin{lstlisting}
\hvFloat[doubleFULLPAGE]%
  {figure}%
  {\includegraphics[height=\paperheight]{images/rheinsberg}}%
  [A doublepage image with a caption on the right side of the right part.]%
  {A caption for a double-sided image that will be placed on the right side of the
   right-hand part of the illustration. The illustration begins on the left edge of 
   the paper. A short form is used for the LOF. 
   The parameter is \texttt{doubleFULLPAGE}}%
  {fig:doubleFULLPAGE0a}
\end{lstlisting}


\marginnote{Fig. \ref{fig:doubleFULLPAGE0a}}
\hvFloat[doubleFULLPAGE]%
  {figure}%
  {\includegraphics[height=\paperheight]{images/rheinsberg}}%
  [A doublepage image with a caption on the right side of the right part.]%
  {A caption for a double-sided image that will be placed on the right side of the
   right-hand part of the illustration. The illustration begins on the left edge of 
   the paper. A short form is used for the LOF. 
   The parameter is \texttt{doubleFULLPAGE}}%
  {fig:doubleFULLPAGE0a}


\Blindtext

\Blindtext

%\Blindtext

\hvblindtext
%\Blindtext

\hvblindtext

\hvblindtext
\hvblindtext

\subsection{The default -- second column}
\begin{lstlisting}
\hvFloat[doubleFULLPAGE]%
  {figure}%
  {\includegraphics[height=\paperheight]{images/rheinsberg}}%
  [A doublepage image with a caption on the right side of the right part.]%
  {A caption for a double-sided image that will be placed on the right side of the
   right-hand part of the illustration. The illustration begins on the left edge of 
   the paper. A short form is used for the LOF. 
   The parameter is \texttt{doubleFULLPAGE}}%
  {fig:doubleFULLPAGE0a2nd}
\end{lstlisting}


\marginnote{Fig. \ref{fig:doubleFULLPAGE0a2nd}}
\hvFloat[doubleFULLPAGE]%
  {figure}%
  {\includegraphics[height=\paperheight]{images/rheinsberg}}%
  [A doublepage image with a caption on the right side of the right part.]%
  {A caption for a double-sided image that will be placed on the right side of the
   right-hand part of the illustration. The illustration begins on the left edge of 
   the paper. A short form is used for the LOF. 
   The parameter is \texttt{doubleFULLPAGE}}%
  {fig:doubleFULLPAGE0a2nd}


\hvblindtext
\hvblindtext

\hvblindtext

\Blindtext

%\Blindtext

%\Blindtext




\subsection{\texttt{bindCorr=1cm} -- first column}

\begin{lstlisting}
\hvFloat[doubleFULLPAGE,capWidth=n,capPos=right,bindCorr=1cm]%
  {figure}%
  {\includegraphics[doubleFULLPAGEbindCorr]{images/rheinsberg}}%
  [A doublepage image with a caption on the right side of the right part.]%
  {A caption for a double-sided image that will be placed on the right side of the
   right-hand part of the illustration. The illustration begins on the left edge of 
   the paper. A short form is used for the LOF. 
   The parameter is \texttt{doubleFULLPAGE}}%
  {fig:doubleFULLPAGE1a}
\end{lstlisting}

\marginnote{Fig. \ref{fig:doubleFULLPAGE1a}}
\hvFloat[doubleFULLPAGE,capWidth=n,capPos=right,bindCorr=1cm]%
  {figure}%
  {\includegraphics[doubleFULLPAGEbindCorr]{images/rheinsberg}}%
  [A doublepage image with a caption on the right side of the right part.]%
  {A caption for a double-sided image that will be placed on the right side of the
   right-hand part of the illustration. The illustration begins on the left edge of 
   the paper. A short form is used for the LOF. 
   The parameter is \texttt{doubleFULLPAGE}}%
  {fig:doubleFULLPAGE1a}

%\hvblindtext

\Blindtext

\Blindtext

%\Blindtext
%\Blindtext

\hvblindtext

\hvblindtext
\hvblindtext

\hvblindtext


\subsection{\texttt{bindCorr=1cm} -- second column}

\begin{lstlisting}
\hvFloat[doubleFULLPAGE,capWidth=n,capPos=right,bindCorr=1cm]%
  {figure}%
  {\includegraphics[doubleFULLPAGEbindCorr]{images/rheinsberg}}%
  [A doublepage image with a caption on the right side of the right part.]%
  {A caption for a double-sided image that will be placed on the right side of the
   right-hand part of the illustration. The illustration begins on the left edge of 
   the paper. A short form is used for the LOF. 
   The parameter is \texttt{doubleFULLPAGE}}%
  {fig:doubleFULLPAGE1a2nd}
\end{lstlisting}

\marginnote{Fig. \ref{fig:doubleFULLPAGE1a2nd}}
\hvFloat[doubleFULLPAGE,capWidth=n,capPos=right,bindCorr=1cm]%
  {figure}%
  {\includegraphics[doubleFULLPAGEbindCorr]{images/rheinsberg}}%
  [A doublepage image with a caption on the right side of the right part.]%
  {A caption for a double-sided image that will be placed on the right side of the
   right-hand part of the illustration. The illustration begins on the left edge of 
   the paper. A short form is used for the LOF. 
   The parameter is \texttt{doubleFULLPAGE}}%
  {fig:doubleFULLPAGE1a2nd}

\hvblindtext

%\Blindtext

\Blindtext

%\Blindtext

\hvblindtext

\hvblindtext




\subsection{\texttt{bindCorr=3mm} -- first column}

\begin{lstlisting}
\hvFloat[doubleFULLPAGE,capWidth=n,capPos=right,bindCorr=3mm]%
  {figure}%
  {\includegraphics[doubleFULLPAGEbindCorr]{images/rheinsberg}}%
  [A doublepage image with a caption on the right side of the right part.]%
  {A caption for a double-sided image that will be placed on the right side of the
   right-hand part of the illustration. The illustration begins on the left edge of 
   the paper. A short form is used for the LOF. 
   The parameter is \texttt{doubleFULLPAGE}}%
  {fig:doubleFULLPAGE2a}
\end{lstlisting}

\marginnote{Fig. \ref{fig:doubleFULLPAGE2a}}
\hvFloat[doubleFULLPAGE,capWidth=n,capPos=right,bindCorr=3mm]%
  {figure}%
  {\includegraphics[doubleFULLPAGEbindCorr]{images/rheinsberg}}%
  [A doublepage image with a caption on the right side of the right part.]%
  {A caption for a double-sided image that will be placed on the right side of the
   right-hand part of the illustration. The illustration begins on the left edge of 
   the paper. A short form is used for the LOF. 
   The parameter is \texttt{doubleFULLPAGE}}%
  {fig:doubleFULLPAGE2a}


\Blindtext

\Blindtext

%\Blindtext
%\Blindtext

\hvblindtext

\hvblindtext

\hvblindtext

\hvblindtext


\subsection{\texttt{bindCorr=3mm} -- second column}

\begin{lstlisting}
\hvFloat[doubleFULLPAGE,capWidth=n,capPos=right,bindCorr=3mm]%
  {figure}%
  {\includegraphics[doubleFULLPAGEbindCorr]{images/rheinsberg}}%
  [A doublepage image with a caption on the right side of the right part.]%
  {A caption for a double-sided image that will be placed on the right side of the
   right-hand part of the illustration. The illustration begins on the left edge of 
   the paper. A short form is used for the LOF. 
   The parameter is \texttt{doubleFULLPAGE}}%
  {fig:doubleFULLPAGE2a2nd}
\end{lstlisting}

\marginnote{Fig. \ref{fig:doubleFULLPAGE2a2nd}}
\hvFloat[doubleFULLPAGE,capWidth=n,capPos=right,bindCorr=3mm]%
  {figure}%
  {\includegraphics[doubleFULLPAGEbindCorr]{images/rheinsberg}}%
  [A doublepage image with a caption on the right side of the right part.]%
  {A caption for a double-sided image that will be placed on the right side of the
   right-hand part of the illustration. The illustration begins on the left edge of 
   the paper. A short form is used for the LOF. 
   The parameter is \texttt{doubleFULLPAGE}}%
  {fig:doubleFULLPAGE2a2nd}


\Blindtext

%\Blindtext

%\Blindtext

%\hvblindtext

\hvblindtext

\hvblindtext

\hvblindtext



\subsection{\texttt{bindCorr=<inside textwidth>} -- first column}

\begin{lstlisting}
\hvFloat[doubleFULLPAGE,capWidth=n,bindCorr=\the\dimexpr1in+\oddsidemargin]%
  {figure}%
  {\includegraphics[doubleFULLPAGEbindCorr]{images/rheinsberg}}%
  [A doublepage image with a caption on the right side of the right part.]%
  {A caption for a double-sided image that will be placed on the right side of the
   right-hand part of the illustration. The illustration begins on the left edge of 
   the paper. A short form is used for the LOF. 
   The parameter is \texttt{doubleFULLPAGE}}%
  {fig:doubleFULLPAGE3a}
\end{lstlisting}

\marginnote{Fig. \ref{fig:doubleFULLPAGE3a}}
\hvFloat[doubleFULLPAGE,capWidth=n,bindCorr=\the\dimexpr1in+\oddsidemargin]%
  {figure}%
  {\includegraphics[doubleFULLPAGEbindCorr]{images/rheinsberg}}%
  [A doublepage image with a caption on the right side of the right part.]%
  {A caption for a double-sided image that will be placed on the right side of the
   right-hand part of the illustration. The illustration begins on the left edge of 
   the paper. A short form is used for the LOF. 
   The parameter is \texttt{doubleFULLPAGE}}%
  {fig:doubleFULLPAGE3a}


\Blindtext

\Blindtext

%\Blindtext

\hvblindtext

\hvblindtext

%\hvblindtext

%\hvblindtext


\subsection{\texttt{bindCorr=<inside textwidth>} -- second column}

\begin{lstlisting}
\hvFloat[doubleFULLPAGE,capWidth=n,bindCorr=\the\dimexpr1in+\oddsidemargin]%
  {figure}%
  {\includegraphics[doubleFULLPAGEbindCorr]{images/rheinsberg}}%
  [A doublepage image with a caption on the right side of the right part.]%
  {A caption for a double-sided image that will be placed on the right side of the
   right-hand part of the illustration. The illustration begins on the left edge of 
   the paper. A short form is used for the LOF. 
   The parameter is \texttt{doubleFULLPAGE}}%
  {fig:doubleFULLPAGE3a2nd}
\end{lstlisting}

\marginnote{Fig. \ref{fig:doubleFULLPAGE3a2nd}}
\hvFloat[doubleFULLPAGE,capWidth=n,bindCorr=\the\dimexpr1in+\oddsidemargin,capPos=after]%
  {figure}%
  {\includegraphics[doubleFULLPAGEbindCorr]{images/rheinsberg}}%
  [A doublepage image with a caption on the right side of the right part.]%
  {A caption for a double-sided image that will be placed on the right side of the
   right-hand part of the illustration. The illustration begins on the left edge of 
   the paper. A short form is used for the LOF. 
   The parameter is \texttt{doubleFULLPAGE}}%
  {fig:doubleFULLPAGE3a2nd}


\Blindtext

\Blindtext

\Blindtext

\hvblindtext



%\cleardoublepage


\hvblindtext

\hvblindtext

\Blindtext

\hvblindtext

\section{Full height, but no complete double page}
\subsection{The default -- in the first column}

\begin{lstlisting}
\hvFloat[doubleFULLPAGE,capPos=right]%
  {figure}%
  {\includegraphics[height=\paperheight]{images/rheinsberg}}%
  [A doublepage image with a caption on the right side of the right part.]%
  {A caption for a double-sided image that will be placed on the right side of the
   right-hand part of the illustration. The illustration begins on the left edge of 
   the paper. A short form is used for the LOF. 
   The parameter is \texttt{doubleFULLPAGE}}%
  {fig:doubleFULLPAGE0n}
\end{lstlisting}


\marginnote{Fig. \ref{fig:doubleFULLPAGE0n}}
\hvFloat[doubleFULLPAGE,capPos=right]%
  {figure}%
  {\includegraphics[height=\paperheight]{images/rheinsberg}}%
  [A doublepage image with a caption on the right side of the right part.]%
  {A caption for a double-sided image that will be placed on the right side of the
   right-hand part of the illustration. The illustration begins on the left edge of 
   the paper. A short form is used for the LOF. 
   The parameter is \texttt{doubleFULLPAGE}}%
  {fig:doubleFULLPAGE0n}

\Blindtext

\Blindtext

%\Blindtext

\clearpage
\Blindtext
%\hvblindtext

%\hvblindtext

\Blindtext

\hvblindtext
\Blindtext

\subsection{The default -- in the second column}


\begin{lstlisting}
\hvFloat[doubleFULLPAGE,capPos=before]%
  {figure}%
  {\includegraphics[height=\paperheight]{images/rheinsberg}}%
  [A doublepage image with a caption on the right side of the right part.]%
  {A caption for a double-sided image that will be placed on the right side of the
   right-hand part of the illustration. The illustration begins on the left edge of 
   the paper. A short form is used for the LOF. 
   The parameter is \texttt{doubleFULLPAGE}}%
  {fig:doubleFULLPAGE02ndn}
\end{lstlisting}


\marginnote{Fig. \ref{fig:doubleFULLPAGE02ndn}}
\hvFloat[doubleFULLPAGE,capPos=before]
  {figure}%
  {\includegraphics[height=\paperheight]{images/rheinsberg}}%
  [A doublepage image with a caption on the right side of the right part.]%
  {A caption for a double-sided image that will be placed on the right side of the
   right-hand part of the illustration. The illustration begins on the left edge of 
   the paper. A short form is used for the LOF. 
   The parameter is \texttt{doubleFULLPAGE}}%
  {fig:doubleFULLPAGE02ndn}

\Blindtext

\Blindtext

\Blindtext

%\hvblindtext

\Blindtext

%\Blindtext


\subsection{The default -- in the second column}


\begin{lstlisting}
\hvFloat[doubleFULLPAGE,capPos=after,twoColumnCaption]%
  {figure}%
  {\includegraphics[doubleFULLPAGE]{images/rheinsberg}}%
  [A doublepage image with a caption on the right side of the right part.]%
  {A caption for a double-sided image that will be placed on the right side of the
   right-hand part of the illustration. The illustration begins on the left edge of 
   the paper. A short form is used for the LOF. 
   The parameter is \texttt{doubleFULLPAGE}}%
  {fig:doubleFULLPAGE02ndnn}
\end{lstlisting}


\marginnote{Fig. \ref{fig:doubleFULLPAGE02ndnn}}
\hvFloat[doubleFULLPAGE,capPos=after,twoColumnCaption]%
  {figure}%
  {\includegraphics[doubleFULLPAGE]{images/rheinsberg}}%
  [A doublepage image with a caption on the right side of the right part.]%
  {A caption for a double-sided image that will be placed 
   at the bottom of the following page. The doublepage image begins on the left edge of 
   the paper. A short form is used for the LOF. 
   The parameter is \texttt{doubleFULLPAGE}}%
  {fig:doubleFULLPAGE02ndnn}

\Blindtext

\Blindtext

\Blindtext

\hvblindtext

\hvblindtext
\hvblindtext



\end{document}